%!TEX TS-program = lualatex
\documentclass[]{friggeri-cv}
\usepackage[none]{hyphenat}
\usepackage{fontspec}
\usepackage{fontawesome}

\begin{document}
\header{Alexandre }{Arpin}
       {Développeur Logiciel}
       {\faCloud ~alexandrearpin.com~ | \faHome ~Montréal, Qc~ | ~ \faPhone ~(514) 730-6126~ | ~\faEnvelopeAlt  ~arpin.alexandre@gmail.com}

\section{éducation}

\begin{entrylist}
  \entry
    {2011-2016}
    {B.Ing {\normalfont Génie logiciel}}
    {École de technologie supérieure}
    {}
    {0cm}
  \entry
    {2008-2011}
    {DEC. {\normalfont Technique d'informatique de gestion}}
    {Collège Ahuntsic}
    {}
    {0cm}
\end{entrylist}

\section{expériences}

\begin{entrylist}
  \entry
    {2016}
    {GSOFT, Montréal}
    {C\#, WPF, ASP.NET MVC, html5, sass, gulp}
    {\emph{Conception, développement et maintenance de Sharegate, une application pour la gouvernance et la migration d'environnements SharePoint.}
    \vspace{0.1cm}
    \begin{itemize}
      \item Développer des fonctionnalités de qualité dans un contexte "Full Stack"
      \item Maintenir des standards de hautes qualités à l'aide de tests unitaires et de documentation
      \item Développer et déployer des applications modulaires à l'aide de Microsoft Azure Services
    \end{itemize}}
    {0.3cm}
  \entry
    {2015}
    {Resulto, Montréal}
    {python, django, django-rest-framework, html5, sass, gulp}
    {\emph{Conception, développement et maintenance de LoyalAction, une solution pour la fidélisation de clientèle.}
    \vspace{0.1cm}
    \begin{itemize}
      \item Développer des fonctionnalités de qualité dans un contexte "Full Stack"
      \item Intégrer des applications et services de tierces parties
      \item Maintenir des standards de hautes qualités à l'aide de tests unitaires et de documentation
      \item Rechercher et implémenter un algorithme de génération de nombre non-séquentiel unique
      \item Publier des outils de développement open-source à l'aide de package PyPi
    \end{itemize}}
    {0.3cm}
  \entry
    {2013-2014}
    {Beehivr Technologies, Repentigny}
    {C\#, ASP.NET MVC, WebApi, html5, css3, Android, Obj-C, JavaFX}
    {\emph{Conception, développement et maintenance de Beehivr, une plateforme de point d'engagement pour iOS et Android.}
    \vspace{0.1cm}
    \begin{itemize}
      \item Développer un tableau de bord d'administration pour gérer des centaines d'appareils iOS
      \item Développer un Api REST pour les applications Beehivr et les fournisseurs
      \item Intégrer Amazon Web Services aux applications et à l'environnement de développement
      \item Collaborer à la conception et aux revues de code de l'application iOS Beehivr
      \item Porter et adapter l'application iOS Beehivr vers Android
    \end{itemize}}
    {0cm}
\end{entrylist}

\section{projets}

\begin{entrylist}
 \entry
    {2016}
    {Root {\normalfont root.plants.ninja}}
    {slack, python, django}
    {Une intégration Slack pour améliorer l'espace de travail en y intégrant des plantes. Root vous aidera à choisir les plantes qui vous convient le mieux et vous notifie lorsqu'elles ont besoin de maintenance.}
    {0.1cm}
 \entry
    {2015}
    {Magic: the Gathering Card Identifier {\normalfont alexandrearpin.com/mtg-card-identifier}}
    {python, django}
    {Librairie et application web permettant d'identifier des cartes Magic: the Gathering prisent en photo à l'aide de vision numérique (OpenCV et Tesseract).}
    {0.1cm}
  \entry
    {2015}
    {Magic: the Gathering Font {\normalfont alexandrearpin.com/mtg-font}}
    {html5, css3}
    {Inspiré par FontAwesome, Web font facile à utiliser pour tous les symboles du jeu Magic: The Gathering.}
    {0.1cm}
  \entry
    {2014}
    {Unfinished Asteroids libGDX {\normalfont alexandrearpin.com/unfinished-asteroids-libgdx}}
    {Android}
    {À des fins d’apprentissage, on vous met dans la peau d'un nouveau développeur se joignant à un projet incomplet du jeu Asteroid. Par la suite, une mise en scène vous guidera afin de résoudre des bogues et ajouter de nouvelles fonctionnalités à ce jeu d’arcade classique.}
    {0cm}
\end{entrylist}
\end{document}
